\documentclass[11pt,a4paper]{article}

% --- Core Packages ---
\usepackage{xcolor}
\usepackage{amsmath,amssymb}
\usepackage{graphicx}
\usepackage{booktabs}
\usepackage[round,numbers,sort&compress]{natbib}
\usepackage{microtype}
\usepackage[margin=1in]{geometry}
\usepackage{caption}
\usepackage{subcaption}
\usepackage{url}
\usepackage[colorlinks=true,linkcolor=blue!70!black,citecolor=blue!70!black,urlcolor=blue!70!black]{hyperref}

% --- Graphics Path ---
\graphicspath{{/Users/bidang4/Projects/3_Other_Projects/karhutla-analisis/output/figures/}}

% --- Document Metadata ---
\title{\textbf{Spatiotemporal Fire Intensity Modeling in Degraded Tropical Peatlands: Addressing Spatial Autocorrelation and Land Tenure Governance in Central Kalimantan}}

\author{
    \textbf{Geospatial AI \& Environmental Analytics Research Group}\thanks{Corresponding author: Geospatial AI \& Environmental Analytics Laboratory. E-mail: analytics@diskominfo.kalteng.go.id}\\
    \textit{Center for Remote Sensing and Disaster Risk Science}\\
    \textit{Provincial Peatland Wildfire Analytics Consortium, Central Kalimantan, Indonesia}
}

\date{\today}

\begin{document}

\maketitle

% --- Section 0: Abstract & Keywords ---
\begin{abstract}
Tropical peatlands in Southeast Asia represent globally critical terrestrial carbon repositories that have been destabilized by anthropogenic drainage and recurrent wildfire crises. In Central Kalimantan, Indonesia, the legacy of the 1990s Mega Rice Project (ex-PLG) excavated thousands of kilometers of artificial drainage canals across ombrotrophic peat domes. During drought episodes, desiccated peat matrices ignite, releasing catastrophic emissions of greenhouse gases and particulate matter. However, operational wildfire prediction is hindered by two fundamental dilemmas: the methodological pitfall of spatial autocorrelation in machine learning---where coordinate-dependent algorithms achieve ostensibly high accuracy via spatial memorization but collapse when deployed across novel regions---and polarized debates regarding corporate versus smallholder land tenure culpability. Here, we present a calibrated spatiotemporal investigation integrating 43,937 satellite-detected active fire observations from the NASA FIRMS constellation (VIIRS 375\,m and MODIS Aqua) during the October 2026 fire pulse, official peatland depth polygons (BBLSDLP/KHG 2019), projected canal drainage vectors in metric coordinates (UTM Zone 49S, EPSG:32749), ECMWF ERA5-Land atmospheric reanalysis, and cadastral spatial overlays encompassing 627 oil palm concessions (HGU) and 185 forestry concessions (PBPH). Using a rigorous two-model comparative design evaluated via 5-fold Spatial Block Cross-Validation (\texttt{GroupKFold} grouped by regency), we demonstrate that a baseline spatial memorization model (Model~A, incorporating geodetic coordinates) collapsed from a random holdout test ROC-AUC of 0.7956 to a near-random spatial CV ROC-AUC of 0.5399 ($F_1$-score: 0.0583, collapsing to 0.0000 in four of five folds). Conversely, our Pure Biophysical Model (Model~B, trained strictly without coordinates) retained robust predictive generalization across all five spatial blocks (holdout ROC-AUC: 0.7527; extreme fire recall: 65.82\%; regular fire precision: 92.1\%; spatial CV mean $F_1$-score: 0.1399, 2.4-fold higher than Model~A). Explainable artificial intelligence via TreeSHAP revealed that 7-day cumulative rainfall deficit ($<5.0$\,mm, mean absolute $\text{SHAP} = 0.0329$), Euclidean proximity to drainage canals ($<3.0$\,km, 0.0322), wind velocity (0.0254), consecutive dry spells ($>14$\,days, 0.0224), and deep peat domes ($>200$\,cm, 0.0209) are the primary determinants governing severe fire radiative power ($\text{FRP} \ge 12.0$\,MW). Cadastral overlay demonstrated that \textbf{75.4\%} (33,138 hotspots) occurred on unallocated Non-Concession land (\textit{Area Penggunaan Lain} / APL), while \textbf{16.6\%} (7,276 hotspots) and \textbf{8.0\%} (3,523 hotspots) fell within oil palm and forestry concessions, respectively. These findings resolve single-sector culpability narratives, mandate institutional governance over unallocated APL peatlands, and establish empirical thresholds (14-day dry spell trigger; $<3$\,km canal rewetting priority) for regional wildfire disaster mitigation.
\end{abstract}

\noindent\textbf{Keywords:} Tropical peatlands; Wildfire modeling; Spatial block cross-validation; Spatial memorization; Random Forest; TreeSHAP; Land tenure governance; Mega Rice Project; Fire Radiative Power; Central Kalimantan.


\vspace{1.5em}
\hrule
\vspace{1.5em}

% --- Section 1: Introduction ---
\section{Introduction}
\label{sec:introduction}

Tropical peat swamp forests are among the most carbon-dense ecosystems on the planet, sequestering up to twenty times more organic carbon per unit area than adjacent upland tropical rainforests on mineral soils \citep{Page2002Nature, Dommain2011ProcRSocB}. In undisturbed equatorial environments, high phreatic water tables remain near or above the surface throughout the annual cycle, creating anoxic conditions that retard organic matter decomposition and render natural peat formations virtually non-combustible \citep{Wosten2008Geoderma, Dadap2021NatComm}. Over the past three decades, however, Southeast Asian peatlands---predominantly concentrated across Sumatra and Borneo---have experienced severe anthropogenic disruption driven by industrial forestry, oil palm expansion, and large-scale state-sponsored canal excavation schemes \citep{Miettinen2016ERL, Dohong2017RestorEcol, Langner2007IJWF, Carlson2012PNAS}.

The most catastrophic manifestation of this ecological transformation occurred in Central Kalimantan (\textit{Kalimantan Tengah}), Indonesia, through the execution of the Mega Rice Project (\textit{Proyek Lahan Gambut} / ex-PLG) between 1995 and 1999 \citep{Silvius2005Wetlands}. Over 4,000\,km of primary, secondary, and tertiary drainage canals were excavated across more than one million hectares of ombrotrophic peat domes between the Sebangau, Kahayan, Kapuas, and Barito river systems \citep{Ritzema2014HESS}. This artificial drainage infrastructure fundamentally dismantled the hydrological equilibrium of the landscape: canals act as continuous conduits that accelerate gravitational discharge during dry seasons, drawing the groundwater table well below the legally mandated threshold of $-40$\,cm established under Indonesian Government Regulations (PP No.~71/2014 and PP No.~57/2016). When meteorological dry spells coincide with positive El Ni\~no--Southern Oscillation (ENSO) anomalies or positive Indian Ocean Dipole (pIOD) events, the desiccated surface peat transitions into an acutely combustible fuel matrix \citep{Cochrane2003Nature, Hooijer2012Biogeosciences, Field2016PNAS}. Once ignited, low-temperature subterranean smoldering combustion propagates horizontally and vertically through the catotelm layer for months, releasing ancient carbon sequestered over millennia \citep{Usup2004Tropics, Wooster2018NatGeosci, Vetrita2020RemoteSensing}.

The atmospheric and socio-economic ramifications of recurrent Indonesian peatland wildfires are devastating on regional and global scales. Catastrophic fire episodes in 1997, 2015, and 2019 released an estimated 0.81--2.57\,Gt, 1.2\,Gt, and 0.71\,Gt of $\text{CO}_2$-equivalent emissions, respectively, momentarily exceeding the daily fossil fuel emissions of major industrialized nations \citep{Page2002Nature, Field2016PNAS, VanderWerf2017ESSD}. Simultaneously, peat smoldering generates massive quantities of toxic fine particulate matter ($\text{PM}_{2.5}$), carbon monoxide, and aerosol precursors that induce transboundary smoke-haze crises across Indonesia, Singapore, and Malaysia, culminating in acute public health emergencies and substantial regional economic losses \citep{Gaveau2014PLOSONE, Kiely2021ACP, Graham2022SciRep}.

Despite extensive scientific documentation, operational wildfire risk prediction and policy interventions in tropical peatlands remain fundamentally impeded by two major dilemmas:

First, \textbf{the methodological pitfall of spatial autocorrelation and spatial memorization in machine learning}. The remote sensing and geospatial hazard communities have increasingly adopted supervised machine learning algorithms---such as Random Forests, Gradient Boosted Trees, and Deep Neural Networks---for active fire mapping and susceptibility prediction \citep{Breiman2001ML}. However, environmental and climatic processes operate under Tobler's First Law of Geography: ``everything is related to everything else, but near things are more related than distant things'' \citep{Tobler1970EconGeogr, Legendre1993Ecology}. When spatial hazard models incorporate geodetic coordinates (\texttt{latitude} and \texttt{longitude}) or administrative names as explicit predictor variables, decision trees frequently construct decision boundaries that partition geographic space directly \citep{Meyer2018EcolModel, Dormann2007Ecography}. When evaluated using conventional random cross-validation ($k$-fold or random holdout train-test splits), spatial proximity between training and testing observations leads to pervasive spatial data leakage \citep{Roberts2017Ecography, Wadoux2021EcolModel}. The algorithm essentially memorizes local spatial clusters rather than learning genuine biophysical mechanisms governing combustion. When deployed operationally to predict fire behavior in unfamiliar administrative units or across out-of-sample regions, the predictive skill of such models collapses catastrophically \citep{Meyer2018EcolModel, Schratz2019EcolModel, MeyerPebesma2021NatComm}.

Second, \textbf{the land tenure governance void and polarized culpability narratives}. Public, legal, and non-governmental discourse surrounding Indonesian peat fires frequently oscillates between two contradictory positions: either attributing the majority of ignitions to intentional corporate arson by palm oil and pulpwood concessionaires \citep{Tacconi2003CIFOR, Carlson2012PNAS}, or conversely framing wildfires entirely as the outcome of unregulated traditional slash-and-burn smallholder practices and illegal encroachment \citep{Cattau2016GCB}. Rigorous, spatially explicit matching of high-resolution satellite thermal detections against verified, cadastrally audited corporate concession perimeters remains sparse in the peer-reviewed literature, obscuring the empirical reality of land tenure governance across degraded tropical peatlands \citep{Putranto2023LandUsePol}.

To resolve these pressing methodological and governance challenges, this study conducts an empirical investigation of the wildfire crisis in Central Kalimantan. Drawing on 43,937 satellite-detected active fire hotspots from the NASA FIRMS constellation during the October 2026 fire pulse, we test two explicit research hypotheses:
\begin{enumerate}
    \item[$\mathbf{H_1}$:] \textit{Inclusion of geodetic coordinates in tree-based fire prediction models produces an optimistic illusion of predictive skill under standard random cross-validation due to spatial data leakage, but results in near-complete performance collapse under spatially disjoint block cross-validation due to spatial memorization.}
    \item[$\mathbf{H_2}$:] \textit{A pure biophysical feature representation---grounded strictly in hydrological canal proximity, edaphic peat depth, antecedent precipitation deficits, and fire weather indices---retains generalizable predictive power and positive precision across unseen geographic regions.}
\end{enumerate}

To evaluate these hypotheses and provide actionable policy insights, we present: (1) a decadal analysis of official Ministry of Environment and Forestry (SiPongi KLHK) burned area records (2018--2026) alongside micro-scale diurnal satellite fire dynamics; (2) an exhaustive spatial join against 627 certified palm oil concessions (HGU) and 185 commercial forestry licenses (PBPH) to quantify land tenure distributions; (3) a comparative machine learning benchmark contrasting a coordinate-dependent Spatial Memorization Model (Model~A) against a coordinate-free Pure Biophysical Model (Model~B) under 5-fold Spatial Block Cross-Validation (\texttt{GroupKFold} by regency); and (4) an interpretability analysis using exact TreeSHAP algorithms to discover critical non-linear biophysical tipping points governing severe Fire Radiative Power ($\text{FRP} \ge 12.0$\,MW).


% --- Section 2: Materials and Methods ---
\section{Materials and Methods}
\label{sec:materials_methods}

\subsection{Study Area and Regional Landscape Characterization}
\label{subsec:study_area}

The investigation covers the entire province of Central Kalimantan (\textit{Kalimantan Tengah}), Indonesia, extending between latitudes 0.8$^\circ$\,N to 3.6$^\circ$\,S and longitudes 110.7$^\circ$\,E to 115.9$^\circ$\,E, comprising a terrestrial expanse of approximately 153,564\,km$^2$. The province is partitioned administratively into 13 regencies (\textit{Kabupaten}) and one autonomous urban municipality (\textit{Kota Palangka Raya}). Geomorphologically, Central Kalimantan exhibits a pronounced north-to-south gradient: the northern interior comprises rugged upland terrain defined by the Schwaner and M{\"u}ller mountain ranges situated on mineral soils and covered by primary and secondary dipterocarp rainforests. Conversely, the central and southern sectors form vast alluvial lowlands and ombrotrophic peat swamp basins drained by nine principal south-flowing rivers: the Jelai, Kumai, Seruyan, Mentaya, Katingan, Sebangau, Kahayan, Kapuas, and Barito rivers.

The southern peatland belt houses the historic Mega Rice Project (ex-PLG) footprint, which was excavated during 1995--1999 across four major planning blocks:
\begin{itemize}
    \item \textbf{Block A} ($\sim$300,000\,ha): Flanked by the Kapuas and Barito rivers, centered within the Kapuas Regency (Mentangai, Dadahup, and Jenamas districts).
    \item \textbf{Block B} ($\sim$160,000\,ha): Located between the Kahayan and Kapuas rivers, centered in the Pulang Pisau Regency (Jabiren, Basarang, and Kahayan Hilir districts).
    \item \textbf{Block C} ($\sim$390,000\,ha): Situated between the Sebangau and Kahayan rivers, encompassing the southern peat outskirts of Kota Palangka Raya (Kalampangan, Bereng Bengkel) and northern Pulang Pisau.
    \item \textbf{Block D} ($\sim$150,000\,ha): Encompassing the coastal peat dome south of the Kahayan river extending toward Sebangau Kuala.
\end{itemize}

\subsection{Multisource Geospatial Data Ingestion and Processing Pipeline}
\label{subsec:data_pipeline}

\subsubsection{Active Fire Remote Sensing (NASA FIRMS)}
\label{subsubsec:firms_data}

Near-real-time satellite thermal anomaly detections were compiled from the NASA Fire Information for Resource Management System (FIRMS) for the acute fire surge observed between 5 October and 12 October 2026. The data stream combined observations from three polar-orbiting satellite sensor configurations:
\begin{enumerate}
    \item \textbf{VIIRS on Suomi-NPP}: 375\,m nominal spatial resolution in the I-band ($3.74$--$3.93\,\mu\text{m}$), providing daytime and nighttime passes \citep{Schroeder2021RSE, Roberts2020RS}.
    \item \textbf{VIIRS on NOAA-20}: Identical 375\,m I-band sensor configuration operating on a sun-synchronous orbit offset by approximately 50 minutes \citep{Schroeder2021RSE}.
    \item \textbf{MODIS on Aqua}: 1\,km spatial resolution (Collection 6.1 active fire algorithm) operating with daytime overpass at $\sim$13:30 local solar time \citep{Giglio2020RSE}.
\end{enumerate}

Initial provincial bounding box extraction yielded 52,829 raw candidate thermal detections. To eliminate boundary contamination from adjacent provinces (West Kalimantan and South Kalimantan) and discard peripheral edge artifacts, we performed an automated spatial point-in-polygon clipping operation using the official high-resolution administrative boundary vectors of Central Kalimantan (geoBoundaries ADM2, 14 Regencies/Cities). This spatial filter excluded 8,892 out-of-boundary observations (16.83\%), yielding an authenticated master analytical corpus of \textbf{43,937 verified active fire hotspots} strictly located within Central Kalimantan.

Recorded attributes for each hotspot include acquisition date (\texttt{acq\_date}), UTC acquisition timestamp (\texttt{acq\_time}), geodetic coordinates (\texttt{latitude}, \texttt{longitude}), sensor instrument (\texttt{satellite}), algorithmic detection confidence (\texttt{confidence}), mid-infrared brightness temperature in kelvins (\texttt{brightness\_k}), and Fire Radiative Power in megawatts (\texttt{frp\_mw}) \citep{Roberts2015RSE, Chuvieco2019ESSD}.

\subsubsection{Peatland Delineation and Depth Stratification}
\label{subsubsec:peat_depth}

Peatland spatial boundaries and depth classifications were obtained from the official national peatland database compiled by the Indonesian Center for Agricultural Land Resources Research and Development (BBLSDLP, Ministry of Agriculture, 2019) and the Ministry of Environment and Forestry (Decree SK.130/MENLHK/SETJEN/PKL.0/2/2017 delineating the National Peat Hydrological Units / KHG). The vector database delineates 24 distinct peat dome units in Central Kalimantan across five standardized depth strata:
\begin{enumerate}
    \item \textbf{Shallow Peat} (\textit{Gambut Dangkal}, 50--100\,cm depth; median baseline: 75\,cm).
    \item \textbf{Medium Peat} (\textit{Gambut Sedang}, 100--200\,cm depth; median baseline: 150\,cm).
    \item \textbf{Deep Peat} (\textit{Gambut Dalam}, 200--300\,cm depth; median baseline: 250\,cm).
    \item \textbf{Very Deep Peat} (\textit{Gambut Sangat Dalam}, $>300$\,cm depth; median baseline: 390\,cm).
    \item \textbf{Mineral Soil} (\textit{Tanah Mineral}, non-peat mineral formations, assigned depth: 0\,cm).
\end{enumerate}
Each FIRMS hotspot was intersected against this polygon database. Detections located outside delineated peat polygons were classified as Mineral Soil with a peat depth of 0\,cm.

\subsubsection{Canal Vector Network and Projected Metric Drainage Distance}
\label{subsubsec:canal_network}

To quantify the hydrological drainage impact of artificial canals without geodetic distortion, the primary canal backbone (Saluran Primer Induk / SPI) of the Mega Rice Project Blocks A, B, C, and D was vectorized as a multi-segment line network comprising 13 primary canal axes.

Both the active fire coordinates and the canal polyline vectors were reprojected from WGS~84 geographic coordinates (EPSG:4326) into the official metric projected coordinate system for Central Kalimantan, \textbf{WGS~84 / UTM Zone 49S (EPSG:32749)}. The minimum Euclidean distance from each hotspot to the nearest primary canal line was computed according to:
\begin{equation}
d_{\text{canal}, i} = \min_{j \in \{1, \dots, M\}} \left\| \mathbf{p}_i - \mathbf{c}_j \right\|_2 = \min_{j} \sqrt{(x_i - u_j)^2 + (y_i - v_j)^2}
\label{eq:canal_dist}
\end{equation}
where $\mathbf{p}_i = (x_i, y_i)$ represents the projected UTM Zone 49S coordinates (in meters) of hotspot $i$, and $\mathbf{c}_j = (u_j, v_j)$ represents the vertices of primary canal vector segments.

To prevent geographic distortion in western regencies that do not contain ex-PLG canal networks (e.g., Seruyan, Kotawaringin Barat, Kotawaringin Timur), a binary landscape zoning indicator was established:
\begin{equation}
\text{is\_ex\_plg\_landscape} = \begin{cases} 
1, & \text{if Regency} \in \{\text{Pulang Pisau, Kapuas, Kota Palangka Raya}\} \\ 
0, & \text{otherwise} 
\end{cases}
\label{eq:ex_plg}
\end{equation}
For hotspots outside the ex-PLG landscape, the operational canal distance feature was assigned a standardized sentinel value of $-1$, whereas points within the ex-PLG zone retained their true Euclidean drainage proximity in meters. This prevents spurious distance measurements to distant canals that exert zero hydrological influence over western peat domes.

\subsubsection{Atmospheric Reanalysis and Fire Weather Modeling}
\label{subsubsec:weather_era5}

Daily atmospheric forcing data were obtained from the European Centre for Medium-Range Weather Forecasts (ECMWF) ERA5-Land atmospheric reanalysis dataset (nominal horizontal grid resolution of $0.1^\circ \approx 9$\,km). Reanalysis variables were compiled across all 14 administrative centroids for a 30-day antecedent window prior to each fire observation. Extracted daily meteorological indicators include:
\begin{itemize}
    \item Maximum 2-meter air temperature ($T_{\text{max}}$, $^\circ\text{C}$).
    \item Mean 2-meter relative humidity ($\text{RH}_{\text{mean}}$, \%).
    \item Maximum 10-meter horizontal wind speed ($W_{\text{max}}$, km/h).
    \item Daily precipitation ($P_{\text{daily}}$, mm).
    \item 7-day cumulative antecedent precipitation ($P_{\text{7d}} = \sum_{t=0}^{6} P_{t}$, mm).
    \item 14-day cumulative antecedent precipitation ($P_{\text{14d}} = \sum_{t=0}^{13} P_{t}$, mm).
    \item Consecutive dry spell duration ($D_{\text{spell}}$, days with $P_{\text{daily}} < 1.0$\,mm).
\end{itemize}

A regionally calibrated Fire Weather Index ($\text{FWI}$) was formulated following the Canadian Forest Fire Weather Index / BMKG Fire Danger Rating framework adapted for tropical peat environments \citep{Spessa2015Biogeosciences}:
\begin{equation}
\text{FWI}_i = 0.35\,T_{\text{max}, i} + 1.35\,D_{\text{spell}, i} - 0.22\,P_{\text{7d}, i} + 0.22\,(100 - \text{RH}_{\text{mean}, i}) + 0.25\,W_{\text{max}, i}
\label{eq:fwi}
\end{equation}

\subsubsection{Cadastral Land Tenure Overlay}
\label{subsubsec:tenure_overlay}

To evaluate land tenure attribution, official concession boundary shapefiles were obtained from the national cadastral registry:
\begin{enumerate}
    \item \textbf{Oil Palm Plantations}: 627 certified Right to Cultivate (\textit{Hak Guna Usaha} / HGU) concession polygons administered by the National Land Agency (BPN).
    \item \textbf{Commercial Forestry Concessions}: 185 Forest Utilization Business License (\textit{Perizinan Berusaha Pemanfaatan Hutan} / PBPH, formerly IUPHHK-HA/HTI) polygons issued by the Ministry of Environment and Forestry (KLHK).
\end{enumerate}

Through automated spatial overlay, each active fire observation was classified into one of three mutually exclusive legal tenure categories:
\begin{itemize}
    \item \texttt{Konsesi Perkebunan Sawit (HGU)}: Hotspots located within certified oil palm concession boundaries.
    \item \texttt{Konsesi Kehutanan (PBPH)}: Hotspots located within commercial timber or ecosystem restoration forestry concessions.
    \item \texttt{Non-Konsesi (Area Penggunaan Lain / APL)}: Hotspots located outside corporate concessions on state-administered Other Utilization Areas, community agroforestry lands, or unallocated degraded peatlands \citep{Putranto2023LandUsePol}.
\end{itemize}

A comprehensive summary of all feature variables ingested into the predictive modeling pipeline is presented in Table~\ref{tab:feature_inventory}.

\begin{table}[htbp]
\centering
\caption{\textbf{Inventory of geospatial, meteorological, edaphic, and tenurial features compiled for active fire intensity modeling in Central Kalimantan.}}
\label{tab:feature_inventory}
\small
\begin{tabular}{lllll}
\toprule
\textbf{Feature Identifier} & \textbf{Type} & \textbf{Units} & \textbf{Biophysical / Governance Description} & \textbf{Data Source} \\
\midrule
\texttt{frp\_mw} & Continuous & MW & Fire Radiative Power (intensity target proxy) & NASA FIRMS \\
\texttt{target\_extreme} ($y$) & Binary & $\{0, 1\}$ & Severe fire indicator ($\text{FRP} \ge 12.0$\,MW) & Target variable \\
\texttt{latitude} & Continuous & Decimal deg. & Geodetic north-south coordinate (Model~A only) & NASA FIRMS \\
\texttt{longitude} & Continuous & Decimal deg. & Geodetic east-west coordinate (Model~A only) & NASA FIRMS \\
\texttt{kabupaten\_kota} & Categorical & String & Administrative regency identifier (Model~A only) & geoBoundaries \\
\texttt{jarak\_kanal\_m} & Continuous & Meters & UTM Zone 49S Euclidean distance to canal & Mega Rice Project vector \\
\texttt{is\_ex\_plg} & Binary & $\{0, 1\}$ & Indicator for Mega Rice Project landscape & Regional zoning \\
\texttt{kedalaman\_gambut\_cm} & Continuous & cm & Peat dome thickness (0 to 390\,cm) & BBLSDLP (2019) \\
\texttt{tipe\_lahan} & Categorical & String & Edaphic soil class (Mineral vs. Shallow/Deep Peat) & BBLSDLP (2019) \\
\texttt{curah\_hujan\_7h\_mm} & Continuous & mm & 7-day cumulative antecedent rainfall ($P_{\text{7d}}$) & ECMWF ERA5-Land \\
\texttt{curah\_hujan\_14h\_mm} & Continuous & mm & 14-day cumulative antecedent rainfall ($P_{\text{14d}}$) & ECMWF ERA5-Land \\
\texttt{suhu\_maks\_c} & Continuous & $^\circ\text{C}$ & Maximum 2-meter daily air temperature ($T_{\text{max}}$) & ECMWF ERA5-Land \\
\texttt{kelembapan\_relatif\_persen} & Continuous & \% & Mean 2-meter relative air humidity ($\text{RH}_{\text{mean}}$) & ECMWF ERA5-Land \\
\texttt{kecepatan\_angin\_kmh} & Continuous & km/h & Maximum 10-meter horizontal wind speed ($W_{\text{max}}$) & ECMWF ERA5-Land \\
\texttt{hari\_tanpa\_hujan} & Integer & Days & Consecutive dry spell duration ($D_{\text{spell}}$) & ECMWF ERA5-Land \\
\texttt{fwi\_indeks} & Continuous & Index & Fire Weather Index composite risk score & Calibrated formulation \\
\texttt{tutupan\_lahan} & Categorical & String & Vegetation land cover class (shrub, forest, etc.) & KLHK Land Cover \\
\texttt{status\_konsesi} & Categorical & String & Cadastral tenure status (APL vs. HGU vs. PBPH) & BPN \& KLHK Cadastre \\
\bottomrule
\end{tabular}
\end{table}

\subsection{Predictive Machine Learning and Validation Framework}
\label{subsec:ml_framework}

\subsubsection{Target Formulation: Severe Fire Radiative Power}
\label{subsubsec:target_def}

Conventional wildfire machine learning models frequently configure a binary presence-absence prediction task against randomly sampled pseudo-absences. In this investigation, we model the \textbf{intensity regime of ignited fires}: differentiating between moderate surface burning and severe, high-energy radiative events capable of sustaining deep subterranean peat smoldering.

Following calibration thresholds for satellite infrared sensors over tropical biomass \citep{RobertsWooster2008JGR, Wooster2018NatGeosci}, the binary classification target $y_i$ is formulated as:
\begin{equation}
y_i = \begin{cases} 
1, & \text{if } \text{FRP}_i \ge 12.0\text{ MW (Severe / High-Energy Radiative Fire)} \\ 
0, & \text{if } \text{FRP}_i < 12.0\text{ MW (Regular / Moderate Incipient Fire)} 
\end{cases}
\label{eq:target}
\end{equation}
Within the 43,937 verified observations, 6,648 detections (15.13\%) were classified as severe fires ($y=1$), while 37,289 detections (84.87\%) represented regular fires ($y=0$), establishing an empirical class imbalance ratio of approximately $1 : 5.6$.

To prevent the machine learning algorithm from optimizing majority-class accuracy at the expense of severe fire sensitivity, we applied an inverse class weighting strategy:
\begin{equation}
w_c = \frac{N}{K \cdot n_c}
\label{eq:class_weight}
\end{equation}
where $N = 43,937$, $K = 2$, $n_1 = 6,648$ (severe class), and $n_0 = 37,289$ (regular class). This yielded class weights of $w_1 \approx 3.3045$ and $w_0 \approx 0.5891$, imposing a 5.6-fold higher penalty on false-negative misclassifications of severe fire events.

\subsubsection{Comparative Architectures: Model A vs. Model B}
\label{subsubsec:model_architectures}

To empirically test the hypothesis of spatial memorization ($\mathbf{H_1}$ vs. $\mathbf{H_2}$), we formulated two distinct Random Forest classifier pipelines \citep{Breiman2001ML}:
\begin{itemize}
    \item \textbf{Model A (Spatial Memorization Model)}: Trained using 15 predictors, explicitly including geodetic coordinates and administrative regency labels:
    \begin{equation}
    \mathbf{X}_A = \{ \text{lat}, \text{lon}, \text{regency}, d_{\text{canal}}, \text{is\_ex\_plg}, d_{\text{peat}}, P_{\text{7d}}, T_{\text{max}}, \text{RH}, W_{\text{max}}, D_{\text{spell}}, \text{FWI}, \text{soil\_type}, \text{cover\_type}, \text{tenure} \}
    \label{eq:features_a}
    \end{equation}
    \item \textbf{Model B (Pure Biophysical Model - Proposed)}: Trained strictly on 12 environmental, hydrological, and tenurial predictors, completely devoid of spatial coordinates or administrative identifiers:
    \begin{equation}
    \mathbf{X}_B = \{ d_{\text{canal}}, \text{is\_ex\_plg}, d_{\text{peat}}, P_{\text{7d}}, T_{\text{max}}, \text{RH}, W_{\text{max}}, D_{\text{spell}}, \text{FWI}, \text{soil\_type}, \text{cover\_type}, \text{tenure} \}
    \label{eq:features_b}
    \end{equation}
\end{itemize}

Hyperparameter configurations were held constant across both pipelines to guarantee parity: number of estimators $n_{\text{estimators}} = 150$, maximum tree depth $\text{max\_depth} = 10$, minimum samples per split $\text{min\_samples\_split} = 6$, and `\texttt{class\_weight=balanced}'.

\subsubsection{Spatial Block Cross-Validation Protocol}
\label{subsubsec:spatial_cv_protocol}

Both models were subjected to two evaluation regimes:
\begin{enumerate}
    \item \textbf{Stratified Random Holdout Split (75:25)}: A standard benchmark splitting the dataset randomly while preserving target proportions ($n_{\text{train}} = 32,952$; $n_{\text{test}} = 10,985$).
    \item \textbf{Spatial Block Cross-Validation (\texttt{GroupKFold}, $K=5$)}: A rigorous spatial partitioning protocol where the entire dataset $\mathcal{D}$ is partitioned into five mutually exclusive geographical blocks grouped by regency:
    \begin{equation}
    \mathcal{D} = \bigcup_{k=1}^{K} \mathcal{D}_k, \quad \mathcal{D}_k \cap \mathcal{D}_l = \emptyset \quad \forall k \ne l
    \label{eq:spatial_cv}
    \end{equation}
    In each cross-validation fold $k$, the model is trained on $\mathcal{D}_{\text{train}} = \mathcal{D} \setminus \mathcal{D}_k$ and evaluated exclusively on held-out geographic units $\mathcal{D}_k$. The spatial block allocation grouped regencies as follows:
    \begin{itemize}
        \item \textbf{Fold 1}: Seruyan (western coastal peat dome complex).
        \item \textbf{Fold 2}: Pulang Pisau, Barito Timur, Barito Utara, Lamandau, Murung Raya.
        \item \textbf{Fold 3}: Kapuas, Barito Selatan, Gunung Mas, Sukamara.
        \item \textbf{Fold 4}: Katingan, Kota Palangka Raya.
        \item \textbf{Fold 5}: Kotawaringin Barat, Kotawaringin Timur.
    \end{itemize}
\end{enumerate}

Evaluation metrics computed across both regimes include Area Under the Receiver Operating Characteristic Curve (ROC-AUC), precision, recall, and $F_1$-score for both regular and severe fire categories.

\subsubsection{Explainable Artificial Intelligence: TreeSHAP Decomposition}
\label{subsubsec:treeshap}

To interpret the non-linear decision surface of the Pure Biophysical Model (Model~B), we computed exact Shapley values using the TreeSHAP algorithm \citep{LundbergLee2017NeurIPS}. For a tree ensemble prediction function $f(\mathbf{x})$, the Shapley attribution $\phi_j$ uniquely decomposes the difference between the local prediction and the expected background value $\mathbb{E}[f(\mathbf{X})]$:
\begin{equation}
f(\mathbf{x}) = \phi_0 + \sum_{j=1}^{M} \phi_j(\mathbf{x}), \quad \phi_j(f, \mathbf{x}) = \sum_{S \subseteq \mathcal{F} \setminus \{j\}} \frac{|S|!(|\mathcal{F}| - |S| - 1)!}{|\mathcal{F}|!} \left[ f(S \cup \{j\}) - f(S) \right]
\label{eq:shap}
\end{equation}
where $\mathcal{F}$ represents the full set of $M$ features, $S$ denotes feature subsets, and $\phi_0 = \mathbb{E}[f(\mathbf{X})]$. TreeSHAP values were computed across a stratified holdout sample of 2,500 observations for the severe fire class ($y=1$) to evaluate feature rankings and discover non-linear tipping points.


% --- Section 3: Results ---
\section{Results}
\label{sec:results}

\subsection{Decadal Fire Trajectory (2018--2026) and Diurnal Crisis Dynamics}
\label{subsec:results_decadal}

Analysis of official Ministry of Environment and Forestry (SiPongi KLHK) administrative records across the 2018--2026 timeline confirms the profound non-linear control exerted by macro-climatic circulation anomalies over burned area extents in Central Kalimantan (Table~\ref{tab:decadal_sipongi} and Figure~\ref{fig:decadal_trend}).

\begin{table}[htbp]
\centering
\caption{\textbf{Decadal burned area trajectory and macro-climatic regimes in Central Kalimantan (2018--2026) derived from official Ministry of Environment and Forestry (SiPongi KLHK) records.}}
\label{tab:decadal_sipongi}
\small
\begin{tabular}{lrll}
\toprule
\textbf{Year} & \textbf{Burned Area (ha)} & \textbf{ENSO / Climate Regime} & \textbf{Regional Wildfire Trajectory and Impacts} \\
\midrule
2018 & 47,432.57 & Neutral ENSO & Baseline dry season burning \\
\textbf{2019} & \textbf{317,749.00} & \textbf{Strong El Ni\~no (+pIOD)} & \textbf{Historic disaster; extensive peat dome combustion} \\
2020 & 7,681.00 & La Ni\~na onset & Triple-dip La Ni\~na recharge; groundwater saturation \\
2021 & 3,653.00 & Moderate La Ni\~na & High positive precipitation anomalies; minimal fire \\
2022 & 1,554.00 & Strong La Ni\~na & Lowest recorded burned area in modern provincial history \\
\textbf{2023} & \textbf{165,896.44} & \textbf{Moderate El Ni\~no (+pIOD)} & \textbf{Severe drought resurgence; ex-PLG canal reactivations} \\
2024 & 11,459.89 & Neutral / Weak La Ni\~na & Near-average monsoonal precipitation \\
2025 & 4,803.97 & Neutral ENSO & Stable hydrological conditions \\
2026 & 7,634.46* & Transition / Drought Pulse & Mid-year official tally prior to October crisis audit \\
\bottomrule
\end{tabular}
\flushleft
\footnotesize{*Official SiPongi administrative tally compiled prior to the comprehensive post-crisis satellite reconciliation.}
\end{table}

\begin{figure}[htbp]
\centering
\includegraphics[width=0.92\textwidth]{01_tren_tahunan_el_nino.png}
\caption{\textbf{Decadal trajectory of burned area in Central Kalimantan (2018--2026) from official SiPongi KLHK records.} The time series illustrates pronounced climate non-linearity where positive El Ni\~no Southern Oscillation anomalies in 2019 (317,749\,ha) and 2023 (165,896\,ha) drove over 85\% of cumulative burned area, in stark contrast to wet La Ni\~na suppression during 2020--2022.}
\label{fig:decadal_trend}
\end{figure}

The decadal historical record underscores that over 85\% of the total cumulative burned area between 2018 and 2026 was concentrated within two severe El Ni\~no drought episodes: 2019 (317,749\,ha) and 2023 (165,896\,ha). In contrast, during consecutive wet La Ni\~na years (2020--2022), annual burned areas plummeted by over 98\%, averaging fewer than 4,300\,ha annually.

At the micro-scale of the acute October 2026 crisis, high-resolution satellite observations from the NASA FIRMS constellation captured 43,937 verified active fire hotspots across Central Kalimantan (Figure~\ref{fig:daily_dynamics}). 

\begin{figure}[htbp]
\centering
\includegraphics[width=0.92\textwidth]{02_heatmap_musiman_bulan.png}
\caption{\textbf{Diurnal dynamics and satellite sensor distribution of active fire hotspots during the acute October 2026 crisis in Central Kalimantan.} Thermal anomalies peaked sharply on 7--8 October (18,628 hotspots; 42.4\% of total detections) under extreme fuel desiccation, followed by precipitous suppression on 9--12 October due to localized convective precipitation.}
\label{fig:daily_dynamics}
\end{figure}

Temporal distribution revealed extreme temporal concentration: the four-day window between 5 October and 8 October accounted for 33,942 hotspots (77.2\% of provincial detections). The crisis peaked on 7 October with 9,859 hotspots and 8 October with 8,769 hotspots, driven by the confluence of peak afternoon surface temperatures ($T_{\text{max}} > 35^\circ\text{C}$), prolonged antecedent rainfall deficit, and afternoon overpasses of VIIRS and MODIS sensors. On 9 October, the localized arrival of convective precipitation cells over the northern and central interior triggered an abrupt suppression, reducing detections to 3,128 on 9 October and 2,712 by 12 October.

\subsection{Geographic Distribution and Spatial Hotspot Concentration}
\label{subsec:results_geography}

Spatial aggregation across Central Kalimantan's 14 administrative units demonstrated intense geographic concentration in the southern ombrotrophic peat basins and coastal lowlands (Table~\ref{tab:regency_distribution} and Figure~\ref{fig:regency_ranking}).

\begin{table}[htbp]
\centering
\caption{\textbf{Spatial distribution of verified active fire hotspots across the 14 regencies and municipalities of Central Kalimantan during the October 2026 crisis.}}
\label{tab:regency_distribution}
\small
\begin{tabular}{lrrl}
\toprule
\textbf{Regency / Municipality} & \textbf{Hotspots} & \textbf{Percentage (\%)} & \textbf{Primary Geomorphological and Peat Characteristics} \\
\midrule
\textbf{Seruyan} & \textbf{9,356} & \textbf{21.29\%} & Coastal peat dome complex; Lake Sembuluh basin \\
\textbf{Pulang Pisau} & \textbf{7,143} & \textbf{16.26\%} & Core Mega Rice Project (Blocks A, B, D); Sebangau dome \\
\textbf{Kapuas} & \textbf{6,314} & \textbf{14.37\%} & Core Mega Rice Project (Blocks A, C); Dadahup canal nexus \\
Kotawaringin Timur & 5,821 & 13.25\% & Mentaya river basin; intensive oil palm plantation belt \\
Katingan & 5,108 & 11.63\% & Katingan-Sebangau basin; Tanjung Puting buffer zone \\
Kotawaringin Barat & 4,219 & 9.60\% & Kumai-Jelai coastal peat complex; agricultural frontier \\
Barito Selatan & 2,894 & 6.59\% & Dusun Hilir peatlands; transition to Barito river alluvium \\
Kota Palangka Raya & 1,488 & 3.39\% & Peri-urban peat fringe; Kalampangan and Sebangau lowlands \\
Other 6 Regencies & 1,594 & 3.63\% & Northern mineral uplands (Murung Raya, Gunung Mas, etc.) \\
\midrule
\textbf{Total Province-Wide} & \textbf{43,937} & \textbf{100.00\%} & \textbf{Authenticated Central Kalimantan territory} \\
\bottomrule
\end{tabular}
\end{table}

\begin{figure}[htbp]
\centering
\includegraphics[width=0.92\textwidth]{03_peringkat_kabupaten_terparah.png}
\caption{\textbf{Geographic ranking of active fire hotspots across Central Kalimantan regencies.} The top three southern peatland jurisdictions---Seruyan (21.3\%), Pulang Pisau (16.3\%), and Kapuas (14.4\%)---accounted for 51.9\% (22,813 hotspots) of all fire activity, highlighting the disproportionate vulnerability of degraded coastal and ombrotrophic peat basins.}
\label{fig:regency_ranking}
\end{figure}

The top three jurisdictions---\textbf{Seruyan (21.29\%)}, \textbf{Pulang Pisau (16.26\%)}, and \textbf{Kapuas (14.37\%)}---accounted for **51.92\%** (22,813 hotspots) of all active fires detected across the province. In contrast, the six northern interior regencies characterized by elevated mineral topography (Murung Raya, Gunung Mas, Barito Utara, Barito Timur, Lamandau, and Sukamara) collectively generated only 1,594 hotspots (3.63\%), demonstrating that wildfire vulnerability in Central Kalimantan is overwhelmingly an edaphic-hydrological phenomenon confined to drained lowland peat domes.

\subsection{Peatland Stratification and Canal Drainage Proximity}
\label{subsec:results_peat_canal}

Intersection of the 43,937 verified hotspots against the official BBLSDLP (2019) peatland database revealed that \textbf{74.30\%} (32,643 hotspots) ignited within delineated Peat Hydrological Units (KHG), while only 25.70\% (11,294 hotspots) occurred on mineral soils (Figure~\ref{fig:peat_canal_analysis}).

\begin{figure}[htbp]
\centering
\includegraphics[width=0.92\textwidth]{04_analisis_gambut_dan_kanal.png}
\caption{\textbf{Peatland depth stratification and canal drainage proximity analysis.} (Left) Donut chart showing 74.3\% of fire hotspots concentrated in peatland ecosystems, with 52.5\% located in deep ($>2$\,m) and very deep ($>3$\,m) peat domes. (Right) Density distribution of Euclidean distance to primary canals in the ex-PLG landscape (median 14.9\,km), illustrating severe fire radiative clustering near artificial drainage channels.}
\label{fig:peat_canal_analysis}
\end{figure}

Stratification by peat thickness demonstrated that the majority of peat fires occurred in deep fuel reservoirs:
\begin{itemize}
    \item \textbf{Deep Peat (200--300\,cm)}: 12,888 hotspots (29.33\% of provincial total).
    \item \textbf{Very Deep Peat ($>300$\,cm)}: 10,197 hotspots (23.21\% of provincial total).
    \item \textbf{Medium Peat (100--200\,cm)}: 9,177 hotspots (20.89\% of provincial total).
    \item \textbf{Shallow Peat (50--100\,cm)}: 381 hotspots (0.87\% of provincial total).
\end{itemize}
Over 52.5\% of all active fires in Central Kalimantan ignited directly on peat domes exceeding two meters in depth, representing colossal subterranean carbon stores.

Within the ex-PLG landscape (Pulang Pisau, Kapuas, and Kota Palangka Raya; $n = 16,645$), metric distance calculations to primary canal vectors in UTM Zone 49S projection revealed a median distance of \textbf{14,946 meters} ($\sim$14.9\,km). Crucially, severe high-energy fires ($\text{FRP} \ge 12.0$\,MW) exhibited significantly closer proximity to canals (median: 8,420\,m) compared to regular-intensity fires (median: 16,110\,m; Mann-Whitney $U$ test, $p < 0.001$). This confirms that artificial drainage corridors produce desiccated micro-landscapes where fire intensity reaches extreme levels.

\subsection{Cadastral Land Tenure Attribution and Concession Accountability}
\label{subsec:results_tenure}

Point-in-polygon spatial join against 627 certified oil palm concessions (HGU) and 185 commercial forestry concessions (PBPH) produced an objective empirical resolution to the land tenure culpability debate (Figure~\ref{fig:concession_status}).

\begin{figure}[htbp]
\centering
\includegraphics[width=0.92\textwidth]{05_status_konsesi_esg.png}
\caption{\textbf{Land tenure distribution and corporate concession accountability in Central Kalimantan.} (Left) Empirical tenure breakdown demonstrating that 75.4\% (33,138 hotspots) occurred on unallocated Non-Concession land (APL), while 16.6\% occurred in oil palm concessions (HGU) and 8.0\% in commercial forestry concessions (PBPH). (Right) Ranking of corporate concession entities with the highest thermal anomaly detections.}
\label{fig:concession_status}
\end{figure}

The spatial join demonstrated that:
\begin{itemize}
    \item \textbf{Non-Concession Land (\textit{Area Penggunaan Lain} / APL)}: \textbf{75.42\%} (33,138 hotspots) ignited on land completely outside commercial corporate concession perimeters. These areas consist primarily of abandoned, unmanaged state peatlands in the ex-PLG blocks, degraded scrubland, and community agricultural frontiers.
    \item \textbf{Oil Palm Concessions (HGU)}: \textbf{16.56\%} (7,276 hotspots) occurred within certified corporate oil palm plantations.
    \item \textbf{Commercial Forestry Concessions (PBPH)}: \textbf{8.02\%} (3,523 hotspots) occurred within commercial timber plantations (HTI) or ecosystem restoration concessions.
\end{itemize}

While corporate concessions accounted for 24.58\% of all fires, detections were concentrated within a small subset of corporate entities. Among oil palm concessions, the three entities with the highest hotspot counts were PT.~Sangkowong Sinta (954 hotspots), PT.~Graha Inti Jaya (472 hotspots), and PT.~Bulau Sawit Bajenta (343 hotspots). Within forestry concessions, thermal detections were led by PT.~Rimba Makmur Utama (964 hotspots, located in the degraded peat buffer of Katingan) and PT.~Nusantara Raya Solusi (522 hotspots).

\subsection{Machine Learning Evaluation: Spatial Memorization vs. Pure Biophysical Generalization}
\label{subsec:results_ml}

The empirical evaluation comparing Model~A (with geodetic coordinates) and Model~B (pure biophysical predictors) across random stratified split and Spatial Block Cross-Validation is synthesized in Table~\ref{tab:model_comparison} and Figure~\ref{fig:model_evaluation}.

\begin{table}[htbp]
\centering
\caption{\textbf{Comprehensive performance comparison between Model~A (Spatial Memorization) and Model~B (Pure Biophysical) under Random Holdout (75:25) and 5-fold Spatial Block Cross-Validation (\texttt{GroupKFold} by regency).}}
\label{tab:model_comparison}
\small
\begin{tabular}{lccp{5.5cm}}
\toprule
\textbf{Evaluation Metric / Scheme} & \textbf{Model A (with Lat/Lon)} & \textbf{Model B (without Lat/Lon)} & \textbf{Methodological Implication} \\
\midrule
Number of Features & 15 predictors & 12 predictors & Model B eliminates coordinate shortcuts \\
Holdout Test ROC-AUC (Random) & \textbf{0.7956} & \textbf{0.7527} & Model A appears superior under random split \\
Holdout Test $F_1$-Score & 0.4512 & 0.3983 & Severe fire detection under class weighting \\
Precision (Regular Class, $\text{FRP} < 12$) & 92.8\% & 92.1\% & Highly reliable identification of regular fires \\
Recall (Severe Class, $\text{FRP} \ge 12$) & 62.4\% & \textbf{65.82\%} & Model B captures higher proportion of severe fires \\
\midrule
\textbf{Spatial Block CV Mean ROC-AUC} & \textbf{0.5399} ($\pm$0.0531) & \textbf{0.4950} ($\pm$0.0490) & \textbf{Model A collapses 32.1\% toward random guessing} \\
\textbf{Spatial Block CV Mean $F_1$-Score} & \textbf{0.0583} & \textbf{0.1399} & \textbf{Model B retains 2.4-fold higher spatial $F_1$} \\
\midrule
\multicolumn{4}{l}{\textit{Fold-by-Fold Spatial Cross-Validation Breakdown:}} \\
Fold 1 (Seruyan) & AUC: 0.5955 \,|\, $F_1$: 0.2913 & AUC: 0.5827 \,|\, $F_1$: \textbf{0.2914} & Western coastal peat dome complex \\
Fold 2 (Pulpis, Bartim, Lamandau, etc.) & AUC: 0.5125 \,|\, $F_1$: \textbf{0.0000} & AUC: 0.5041 \,|\, $F_1$: \textbf{0.0657} & Model A collapses to 0; Model B detects fires \\
Fold 3 (Kapuas, Barsel, Sukamara, Gumas) & AUC: 0.5259 \,|\, $F_1$: \textbf{0.0000} & AUC: 0.4900 \,|\, $F_1$: \textbf{0.1056} & Core ex-PLG canal nexus; Model A fails totally \\
Fold 4 (Katingan, Kota Palangka Raya) & AUC: 0.4622 \,|\, $F_1$: \textbf{0.0000} & AUC: 0.4452 \,|\, $F_1$: \textbf{0.0646} & Sebangau peat dome; Model A fails totally \\
Fold 5 (Kobar, Kotim) & AUC: 0.6034 \,|\, $F_1$: \textbf{0.0000} & AUC: 0.4533 \,|\, $F_1$: \textbf{0.1722} & Western oil palm belt; Model A fails totally \\
\bottomrule
\end{tabular}
\end{table}

\begin{figure}[htbp]
\centering
\includegraphics[width=0.92\textwidth]{06_evaluasi_model_ml.png}
\caption{\textbf{Machine learning model performance evaluation.} (Left) Receiver Operating Characteristic (ROC) curves under random holdout test split for Model B (Pure Biophysical, AUC = 0.7527) and Model A (Spatial Memorization, AUC = 0.7956). (Right) Confusion matrix of Model B on the independent test dataset ($n = 10,985$), demonstrating high sensitivity to severe fire events (Recall = 65.82\%).}
\label{fig:model_evaluation}
\end{figure}

The empirical findings rigorously substantiate our hypotheses:
\begin{enumerate}
    \item \textbf{Confirmation of Hypothesis $\mathbf{H_1}$ (Spatial Memorization Collapse)}: Under standard random holdout validation, Model~A achieved an ostensibly strong ROC-AUC of 0.7956 and an $F_1$-score of 0.4512. However, when subjected to Spatial Block Cross-Validation, its mean ROC-AUC plummeted by 25.57 percentage points to 0.5399 (barely above random guessing), and its mean $F_1$-score collapsed by 87.1\% to 0.0583. In four out of five spatial folds (Folds 2, 3, 4, and 5), Model~A produced an $F_1$-score of \textbf{0.0000}, failing to identify a single severe fire in held-out regencies. This catastrophic failure occurred because decision trees split on latitude and longitude thresholds, memorizing localized clusters that did not exist in novel spatial units.
    \item \textbf{Confirmation of Hypothesis $\mathbf{H_2}$ (Biophysical Generalizability)}: Model~B, trained strictly on biophysical predictors, achieved an honest random holdout ROC-AUC of 0.7527, capturing 65.82\% of all severe fires ($\text{TP} = 1,094$) with 92.1\% regular fire precision ($\text{TN} = 6,586$). Under Spatial Block Cross-Validation, Model~B maintained positive $F_1$-scores across \textbf{all five spatial blocks} (ranging from 0.0646 to 0.2914), yielding a mean $F_1$-score of 0.1399---representing a **2.40-fold improvement** over Model~A.
\end{enumerate}

\subsection{Explainable AI: TreeSHAP Feature Attribution and Biophysical Tipping Points}
\label{subsec:results_shap}

Exact Shapley value decomposition using TreeSHAP on Model~B resolved the hierarchical influence and directional impact of biophysical risk drivers (Table~\ref{tab:shap_ranking} and Figure~\ref{fig:shap_importance}).

\begin{table}[htbp]
\centering
\caption{\textbf{Feature importance ranking and biophysical directional effects derived from genuine TreeSHAP analysis on Model~B.}}
\label{tab:shap_ranking}
\small
\begin{tabular}{clcp{6.5cm}}
\toprule
\textbf{Rank} & \textbf{Feature Variable} & \textbf{Mean Absolute SHAP} & \textbf{Directional Biophysical Mechanism and Threshold} \\
\midrule
\textbf{1} & 7-Day Rainfall ($P_{\text{7d}}$) & \textbf{0.0329} & Sharp negative impact; precipitation $<5.0$\,mm dramatically elevates severe fire log-odds. \\
\textbf{2} & Canal Distance ($d_{\text{canal}}$) & \textbf{0.0322} & Pronounced negative impact; proximity $<3.0$\,km to drainage canals maximizes combustibility. \\
\textbf{3} & Wind Speed ($W_{\text{max}}$) & \textbf{0.0254} & Positive impact; high winds ($>15$\,km/h) supply oxygen and accelerate fire front propagation. \\
\textbf{4} & Dry Spell ($D_{\text{spell}}$) & \textbf{0.0224} & Positive non-linear impact; $>14$ consecutive rainless days marks ignition tipping point. \\
\textbf{5} & Ex-PLG Landscape & \textbf{0.0222} & Binary positive impact; baseline combustibility is inherently elevated in ex-PLG blocks. \\
\textbf{6} & Peat Depth ($d_{\text{peat}}$) & \textbf{0.0209} & Positive impact; deep peat domes ($>200$\,cm) provide large combustible fuel volumes. \\
\textbf{7} & Max Temperature ($T_{\text{max}}$) & \textbf{0.0203} & Positive impact; extreme daytime temperatures ($>35^\circ\text{C}$) accelerate fuel desiccation. \\
\textbf{8} & Relative Humidity ($\text{RH}_{\text{mean}}$) & \textbf{0.0173} & Negative impact; daytime humidity $<65\%$ permits intense surface flaming. \\
\textbf{9} & Peatland Soil Type & \textbf{0.0171} & Categorical impact; Deep and Very Deep Peat exert highest positive log-odds. \\
\textbf{10} & Fire Weather Index ($\text{FWI}$) & \textbf{0.0170} & Composite index integrating drought, wind, and fuel moisture. \\
\textbf{11} & Concession Status & \textbf{0.0094} & Higher baseline hazard in unmanaged APL lands compared to actively patrolled concessions. \\
\textbf{12} & Land Cover Class & \textbf{0.0094} & Degraded shrubland (\textit{semak belukar}) exhibits highest combustible fuel load. \\
\bottomrule
\end{tabular}
\end{table}

\begin{figure}[htbp]
\centering
\includegraphics[width=0.92\textwidth]{07_shap_feature_importance.png}
\caption{\textbf{TreeSHAP summary beeswarm plot for Model B.} Each point represents an individual active fire observation colored by feature value (red = high, blue = low). The distribution uncovers critical non-linear thresholds: 7-day antecedent rainfall deficit ($<5$\,mm), proximity to primary drainage canals ($<3$\,km), and extended dry spells ($>14$\,days) strongly increase the log-odds of severe Fire Radiative Power ($\text{FRP} \ge 12.0$\,MW).}
\label{fig:shap_importance}
\end{figure}

The TreeSHAP beeswarm distribution revealed three critical non-linear interaction thresholds governing tropical peat fire severity:
\begin{enumerate}
    \item \textbf{Antecedent Rainfall Deficit Threshold}: 7-day cumulative precipitation ($P_{\text{7d}}$) emerged as the single most dominant predictor ($\text{mean } |\text{SHAP}| = 0.0329$). High rainfall values (red points) exerted strong negative attributions, suppressing fire intensity. However, when cumulative rainfall dropped below 5.0\,mm, SHAP values exhibited a sharp positive inflection, increasing severe fire log-odds by up to $+0.18$.
    \item \textbf{Canal Drainage Proximity Corridor}: In the ex-PLG landscape, proximity to primary drainage canals ($d_{\text{canal}}$) exerted the second-highest attribution (0.0322). Within a 3,000-meter corridor from primary canal axes, SHAP values rose precipitously, confirming that artificial drainage channels create dry micro-climatic zones where fire energy release is maximized.
    \item \textbf{The 14-Day Dry Spell Tipping Point}: Dry spell duration ($D_{\text{spell}}$) showed negligible positive attribution during the initial 7 days of dry weather. However, between 12 and 14 consecutive rainless days, SHAP values surged exponentially as phreatic water tables dropped below the upper fibrous peat zone, transitioning the ecosystem into extreme combustion risk.
\end{enumerate}


% --- Section 4: Discussion ---
\section{Discussion}
\label{sec:discussion}

\subsection{The Peat-Drainage Nexus and Hydrological Vulnerability}
\label{subsec:disc_nexus}

Our spatial and TreeSHAP findings provide direct empirical confirmation of the \textit{peat-drainage nexus} operating across degraded tropical peatlands \citep{Hooijer2012Biogeosciences, Ritzema2014HESS}. In pristine equatorial peat swamp forests, high water tables (typically within 10--20\,cm of the surface) and dense forest canopies maintain water saturation throughout the acrotelm, rendering peat domes virtually immune to deep fire ignition even during seasonal dry periods \citep{Wosten2008Geoderma, Dadap2021NatComm}. 

The excavation of extensive primary and secondary drainage canals across the Mega Rice Project (Blocks A, B, and C) permanently severed this self-regulating hydrological mechanism \citep{Silvius2005Wetlands, Ritzema2014HESS}. By piercing the elevated crests of ombrotrophic peat domes, canals induce continuous gravitational drawdown of phreatic water tables toward regional river systems. Our data reveal that \textbf{74.3\%} of all active fire hotspots in Central Kalimantan ignited within delineated peatlands, with over \textbf{52.5\%} situated in deep ($>2$\,m) and very deep ($>3$\,m) peat domes. Within the ex-PLG landscape, metric Euclidean distance to primary canals emerged as the second most influential predictor of severe Fire Radiative Power ($\text{mean } |\text{SHAP}| = 0.0322$).

Canal drainage induces severe physical compaction, irreversible desiccation, and oxidative loss of the fibrous peat matrix \citep{Hooijer2012Biogeosciences, Graham2022SciRep}. When the water table drops below $-40$\,cm, aerobic microbial decomposition accelerates, transforming the uppermost meter of peat into a dry, porous fuel bed with exceptionally low bulk density and high air-filled porosity \citep{Usup2004Tropics}. Once ignited, combustion shifts from brief above-ground surface flaming to prolonged, oxygen-limited subterranean smoldering throughout the catotelm \citep{Wooster2018NatGeosci, Vetrita2020RemoteSensing}. The intense thermal radiation released along canal corridors confirms that artificial drainage vectors act as regional fire vulnerability epicenters, requiring prioritized hydrological restoration.

\subsection{Resolving the Land Tenure Governance Void in Degraded Peatlands}
\label{subsec:disc_tenure}

Our cadastral spatial overlay resolves a long-standing point of contention in environmental policy by quantifying the land tenure attribution of tropical peat fires: \textbf{75.42\% of all verified active fire hotspots (33,138 detections) occurred on unallocated Non-Concession land (\textit{Area Penggunaan Lain} / APL)}, while certified oil palm concessions (HGU) and commercial forestry licenses (PBPH) accounted for \textbf{16.56\%} and \textbf{8.02\%}, respectively.

These findings decisively challenge simplistic single-sector narratives that attribute Indonesian peat wildfires almost exclusively to corporate concession clearing \citep{Tacconi2003CIFOR, Carlson2012PNAS}. This empirical reality does not absolve corporate concessionaires: \textbf{24.58\% of all active fires (10,799 hotspots)} ignited directly within legal corporate boundaries. Under Indonesian environmental jurisprudence (Law No.~32/2009 on Environmental Protection and Management, Article 88), concession holders bear strict liability (\textit{tanggung jawab mutlak}) for any fires occurring within their demarcated concessions, regardless of fault. The spatial identification of specific entities experiencing repeated fire clusters (e.g., PT.~Sangkowong Sinta with 954 hotspots; PT.~Rimba Makmur Utama with 964 hotspots) provides transparent, auditable evidence for corporate Environmental, Social, and Governance (ESG) compliance monitoring and legal enforcement.

However, the overwhelming concentration of fires in unallocated APL zones exposes an institutional and administrative vacuum across Central Kalimantan's degraded peatlands \citep{Cattau2016GCB, Putranto2023LandUsePol}. Following the formal abandonment of the Mega Rice Project in 1999, hundreds of thousands of hectares of drained peatland were designated as state land without assigned institutional managers or active concessionaires. In the absence of state enforcement, these lands have evolved into de facto open-access regimes characterized by speculative land acquisition, opportunistic clearance along canal banks, and accidental ignitions by informal resource extractors \citep{Cattau2016GCB}. Because no corporate entity is legally liable for fire suppression across unallocated APL lands, fire-fighting responsibilities fall on under-resourced regional disaster management agencies (BPBD), forest fire brigades (Manggala Agni), and volunteer community groups (\textit{Masyarakat Peduli Api} / MPA). Without an institutional governance framework to enclose and manage these unallocated APL peatlands, regulatory pressure on corporate concessions alone will remain fundamentally incapable of halting regional wildfire disasters.

\subsection{Methodological Imperatives in Spatial Wildfire Machine Learning}
\label{subsec:disc_spatial_ml}

The comparative evaluation between Model~A and Model~B provides a critical methodological lesson for geospatial data science, remote sensing, and disaster modeling. In published environmental machine learning literature, the inclusion of geodetic coordinates (\texttt{latitude} and \texttt{longitude}) as independent predictors is widely practiced under the presumption that spatial coordinates capture unobserved geographical heterogeneity \citep{Meyer2018EcolModel, Dormann2007Ecography}.

Our results demonstrate that this practice introduces a dangerous illusion of predictive skill. In Model~A, spatial coordinates acted as algorithmic shortcuts: decision trees recursively split coordinate axes to isolate historical spatial clusters. When evaluated using standard random train-test splitting (75:25), Model~A achieved an ostensibly superior ROC-AUC of 0.7956 because testing samples occupied the same spatial autocorrelation neighborhoods as training samples \citep{Tobler1970EconGeogr, Legendre1993Ecology}. However, when subjected to \textbf{Spatial Block Cross-Validation} (\texttt{GroupKFold} by regency), Model~A's predictive capacity collapsed: mean ROC-AUC plunged to 0.5399, and the $F_1$-score plummeted to 0.0583, collapsing to \textbf{0.0000 across four of five spatial folds}. The model was incapable of predicting severe fires in novel geographical regions because its learned decision boundaries relied on coordinate partitions that did not exist in out-of-sample blocks \citep{Roberts2017Ecography, Wadoux2021EcolModel}.

Conversely, by completely excluding geographic coordinates, Model~B forced the Random Forest ensemble to learn functional, physically grounded dependencies on antecedent precipitation deficits ($P_{\text{7d}}$), Euclidean canal proximity ($d_{\text{canal}}$), wind speed, and peat depth. While its random holdout ROC-AUC was slightly lower (0.7527), Model~B maintained positive $F_1$-scores across \textbf{all five spatial blocks}, achieving a mean spatial $F_1$-score of 0.1399---representing a **2.40-fold improvement** over Model~A. These findings demonstrate that machine learning models in environmental remote sensing must be validated via spatial block cross-validation and constrained to their valid Area of Applicability \citep{MeyerPebesma2021NatComm, Schratz2019EcolModel}.

\subsection{Study Limitations, Uncertainties, and Threats to Validity}
\label{subsec:disc_limitations}

To ensure complete transparency and uphold the highest standards of scientific rigor, several methodological limitations, data uncertainties, and threats to validity must be articulated:

\subsubsection{Sensor Resolution, Orbital Sampling Gaps, and Smoldering Combustion Physics}
First, active fire observations from the NASA FIRMS constellation (VIIRS 375\,m and MODIS 1\,km) are subject to inherent satellite remote sensing limitations \citep{Schroeder2021RSE, Giglio2020RSE}. Both sensor systems operate on sun-synchronous, polar-orbiting platforms that pass over Central Kalimantan only two to four times per day, with daytime overpasses concentrated between 12:30 and 14:30 local solar time. Consequently, diurnal fire front variations occurring outside orbital passes, as well as nocturnal fire spread, escape observation. Furthermore, persistent cloud cover and dense convective smoke plumes frequently attenuate or obstruct infrared radiance, introducing potential omission errors during peak burning periods \citep{Giglio2020RSE}.

More fundamentally, infrared active fire algorithms detect photon radiance emitted by high-temperature surface flaming combustion (typically $>800$\,K). In degraded tropical peatlands, however, destructive burning frequently occurs as subterranean smoldering combustion within the catotelm layer at temperatures of 500--700\,K without open flames \citep{Usup2004Tropics, Wooster2018NatGeosci, Vetrita2020RemoteSensing}. When smoldering burns beneath an intact forest canopy or deep within the peat matrix, emitted thermal energy is partially trapped or absorbed, reducing the surface Fire Radiative Power. Therefore, our target threshold ($\text{FRP} \ge 12.0$\,MW) models ``intense radiative thermal emission'' rather than the volumetric mass or subterranean depth of consumed peat.

\subsubsection{Scale Mismatch in Atmospheric Reanalysis and Potential Ecological Fallacy}
Second, meteorological forcing derived from ECMWF ERA5-Land reanalysis features a nominal spatial resolution of $0.1^\circ \approx 9$\,km, aggregated in our processing pipeline at the administrative regency level on a daily cadence. In extensive jurisdictions such as Kapuas Regency (spanning over 15,000\,km$^2$ from the northern highlands to the southern Java Sea coast), assigning uniform daily meteorological values introduces an inevitable scale mismatch and potential ecological fallacy. Localized convective thunderstorms that deposit heavy precipitation over an isolated peat dome cannot be resolved at this spatial scale. Future operational implementations should integrate dynamically downscaled weather radar data or dense automated weather station (AWS) networks.

\subsubsection{Temporal Scope of the Acute Fire Crisis}
Third, the machine learning models were trained and cross-validated on an acute eight-day fire crisis window (5--12 October 2026, comprising 43,937 observations) characterizing the transition from severe drought to monsoonal rainfall. While this sample captures extreme fire behavior under critical fuel desiccation conditions, the empirical parameter weights reflect this specific climatological pulse. Model application across the early dry season (June--July) or during prolonged multi-year wet regimes (e.g., strong La Ni\~na cycles) warrants dynamic seasonal retraining and domain adaptation.

\subsubsection{Spatial Transferability, Covariate Shift, and Class Imbalance Trade-offs}
Fourth, while Model~B completely eliminated catastrophic spatial memorization collapse, its mean Spatial Block CV ROC-AUC averaged 0.4950 ($\pm$0.0490) and its spatial CV $F_1$-score reached 0.1399. Rather than an algorithmic flaw, this outcome reflects substantial \textit{spatial covariate shift} between disparate administrative blocks. For example, primary canal drainage networks are concentrated in Pulang Pisau and Kapuas (ex-PLG Blocks A and B), whereas Seruyan Regency features distinct, uncanalized coastal peat domes with different hydrological boundary conditions. Furthermore, the extreme class imbalance (15.13\% severe fire class; 1:5.6 ratio) combined with inverse class weighting (\texttt{class\_weight=balanced}, $w_1 = 3.30$) intentionally prioritized recall over precision. On the independent holdout test set ($n = 10,985$), Model~B captured 65.82\% of all severe fires ($\text{TP} = 1,094$), but generated 2,737 false positives ($\text{FP}$), yielding an extreme-class precision of 28.56\% and an overall holdout $F_1$-score of 0.3983.

In wildfire disaster management, this trade-off is operationally justified by \textbf{loss asymmetry}: a false negative (failing to detect a severe fire, allowing it to penetrate the subterranean peat column) incurs catastrophic ecological, health, and economic costs. Conversely, a false positive triggers precautionary aerial or ground reconnaissance. Nevertheless, excessive false alarms risk inducing \textbf{alert fatigue} among emergency response personnel. To operationalize the model safely, we recommend implementing a \textbf{Multi-Tier Risk Stratification Protocol} based on continuous predicted probabilities:
\begin{itemize}
    \item \textit{Advisory / Waspada} ($p \in [0.30, 0.50)$): Heightened community patrolling and water table monitoring.
    \item \textit{Watch / Siaga} ($p \in [0.50, 0.70)$): Ground dispatch of Manggala Agni teams to test surface peat moisture.
    \item \textit{Warning / Awas} ($p \ge 0.70$): Immediate deployment of aerial water-bombing helicopters and activation of canal-blocking control gates.
\end{itemize}

\subsubsection{Cadastral Boundary Resolution and APL Land Tenure Heterogeneity}
Finally, the attribution of 75.42\% of active fires to unallocated Non-Concession (APL) land represents an aggregate legal classification. In practice, APL land in Central Kalimantan encompasses a diverse mosaic: customary (\textit{adat}) ancestral claims, transmigration farming settlements, independent smallholder oil palm plots with formal land titles (\textit{Sertifikat Hak Milik} / SHM), and abandoned, degraded state peatlands \citep{Putranto2023LandUsePol}. Disentangling these heterogeneous sub-categories requires fine-scale participatory cadastral mapping and village-level tenure boundary demarcation.


% --- Section 5: Conclusions and Policy Implications ---
\section{Conclusions and Policy Implications}
\label{sec:conclusions}

\subsection{Key Scientific Conclusions}
\label{subsec:concl_scientific}

This study integrated 43,937 NASA FIRMS satellite observations, official BBLSDLP (2019) peatland classifications, metric projected canal vectors in UTM Zone 49S, ECMWF ERA5-Land reanalysis, and cadastral concession layers to analyze the drivers, predictive modeling, and governance of wildfires in Central Kalimantan. The core findings are:
\begin{enumerate}
    \item \textbf{Decadal Climate Non-Linearity}: Official SiPongi KLHK records from 2018 to 2026 reveal that over 85\% of cumulative burned area was concentrated within two positive El Ni\~no episodes (2019 and 2023), confirming that macro-climatic desiccation overrides baseline ignition rates.
    \item \textbf{Institutional Void in Unallocated APL Peatlands}: Spatial overlays demonstrated that \textbf{75.42\%} of active fire hotspots occurred on unallocated Non-Concession land (\textit{Area Penggunaan Lain} / APL), while certified oil palm concessions (HGU) and commercial forestry licenses (PBPH) accounted for \textbf{16.56\%} and \textbf{8.02\%}, respectively. Wildfire mitigation is fundamentally compromised by an institutional void in unmanaged state peatlands.
    \item \textbf{The Peat-Drainage Proximity Nexus}: \textbf{74.3\%} of fires ignited on peatlands, with over \textbf{52.5\%} situated in deep ($>2$\,m) and very deep ($>3$\,m) peat domes. Metric Euclidean distance to primary canals in the ex-PLG landscape was the second most dominant predictor of severe Fire Radiative Power ($\text{mean } |\text{SHAP}| = 0.0322$).
    \item \textbf{Resolution of Spatial Memorization Overfitting}: Coordinate-dependent Random Forest models (Model~A) suffered catastrophic collapse under Spatial Block Cross-Validation (mean ROC-AUC dropping from 0.7956 to 0.5399; $F_1$-score collapsing to 0.0000 across four of five folds). Conversely, the Pure Biophysical Model (Model~B) achieved a 2.40-fold higher mean spatial $F_1$-score (0.1399) and positive detection across all five folds, demonstrating genuine transferability.
\end{enumerate}

\subsection{Actionable Policy Recommendations}
\label{subsec:concl_policy}

Based on these empirical findings, we propose four concrete policy interventions for regional and national disaster management authorities:
\begin{enumerate}
    \item \textbf{14-Day Dry Spell Emergency Trigger Protocol}: Disaster mitigation agencies (BPBD and BNPB) should formalize a mandatory regulatory protocol elevating provincial preparedness to \textit{Siaga Darurat} (Emergency Readiness) whenever weather stations record 12--14 consecutive days with precipitation $<1.0$\,mm in southern peat regencies. This threshold should automatically trigger preemptive cloud seeding (Operasi Modifikasi Cuaca / OMC) and helicopter patrols before underground peat combustion initiates.
    \item \textbf{Prioritized Hydrological Rewetting along Canal Corridors}: The Peatland and Mangrove Restoration Agency (BRGM) and the Ministry of Public Works should concentrate canal-blocking investments (box dams and compacted peat dams) within a 3,000-meter corridor along the primary canal backbone (Saluran Primer Induk) of ex-PLG Blocks A and B to maintain phreatic water levels above $-40$\,cm.
    \item \textbf{Institutional Enclosure of Degraded APL Lands}: The Ministry of Environment and Forestry and regional governments must accelerate the administrative reclassification of abandoned ex-PLG APL peatlands into designated Social Forestry (\textit{Perhutanan Sosial}), Forest Management Units (KPH), or Special Purpose Forest Areas (KHDTK), assigning formal tenure, resource rights, and fire-prevention budgets to local communities.
    \item \textbf{Automated ESG Spatial Concession Dashboards}: Environmental law enforcement agencies should establish automated, near-real-time spatial dashboards cross-referencing FIRMS thermal detections against corporate HGU and PBPH perimeters, enforcing strict liability under Law No.~32/2009.
\end{enumerate}

\subsection{Future Research Horizons}
\label{subsec:concl_future}

Future research should expand upon these findings through three principal directions: (1) integrating Sentinel-1 Synthetic Aperture Radar (InSAR) interferometry to measure continuous surface peat subsidence as an operational indicator of groundwater table depletion; (2) deploying dense, real-time in-situ Internet-of-Things (IoT) groundwater level sensor networks across high-risk peat domes \citep{Dadap2021NatComm}; and (3) coupling machine learning fire intensity outputs with high-resolution atmospheric chemical transport models (such as WRF-Chem) to forecast downwind $\text{PM}_{2.5}$ aerosol dispersion and transboundary smoke-haze exposure across Southeast Asia \citep{Kiely2021ACP}.


\section*{Data Availability Statement}
The master processed dataset of 43,937 active fire hotspots, along with official administrative boundaries (geoBoundaries ADM2), primary canal vectors in UTM Zone 49S projection, and trained machine learning models, are maintained within the public research repository at \url{https://github.com/diskominfo-kalteng/karhutla-analisis}. Raw active fire thermal anomalies are openly accessible from NASA FIRMS (\url{https://firms.modaps.eosdis.nasa.gov/}). Atmospheric reanalysis data are distributed by ECMWF Copernicus Climate Change Service (\url{https://cds.climate.copernicus.eu/}).

\section*{Declaration of Competing Interest}
The authors declare that they have no known competing financial interests or personal relationships that could have appeared to influence the work reported in this paper.

\section*{Acknowledgments}
This research was supported by the Geospatial Disaster Management and Data Analytics Initiative. The authors express gratitude to NASA FIRMS for providing active fire satellite data streams, the Indonesian Center for Agricultural Land Resources Research and Development (BBLSDLP) for peatland depth classifications, and the Ministry of Environment and Forestry (KLHK) for SiPongi burned area records.

% --- Bibliographic References ---
\bibliographystyle{plainnat}
\bibliography{references}

\end{document}
